\documentclass[12pt,reqno]{amsart}
\usepackage{/Users/johnmyers/GoogleDrive/tex/handout,/Users/johnmyers/GoogleDrive/tex/customcommands,polynom}
\usepackage{caption}

\hdr{MAT 240}{\S\S19.3, 19.4 The Divergence Theorem}

\begin{document}

\textbf{Suggested problems \S19.3:} 3-9 odd \medskip

\textbf{Suggested problems \S19.4:} 1-19 odd

\bigskip
\hrule

\bigskip

\prob Using the diagrams below, compute the vector area element $d\vec{A}$ for cylindrical and spherical surfaces. Assume the cylinder and sphere have radius $R$.

\bigskip
\begin{minipage}{0.5\textwidth}
\begin{center}
\includegraphics[scale=1.25]{pic7.pdf}
\end{center}
\end{minipage}
\begin{minipage}{0.5\textwidth}
\begin{center}
\includegraphics[scale=1.25]{pic8.pdf}
\end{center}
\end{minipage}

\vspace{1in}
\prob Compute the flux of
	\[\vec{F}(x,y,z) = \ang{x ,y, 3xz^2}
	\]
through the surface $S$ of a cylinder of radius 10, centered on the $z$-axis, oriented outward, with $0 \leq z \leq 5$.

\newpage
\prob Compute the flux of
	\[\vec{F}(x,y,z) = \ang{0,0,z}
	\]
through the sphere $S$ of radius 5, centered at the origin, and oriented outward.

\vfill
\prob Compute the divergence $\nabla \cdot \vec{F}$ for the following vector fields.

\medskip
\begin{enumerate}
\item $\vec{F}(x,y,z) = \ang{-x+y,y+z,-z + x}$\vfill
\item $\vec{F}(x,y,z) = 3x^2 \vec{i} - \sin{(xz)} \vec{j}+3x^2yz^2\vec{k}$
\end{enumerate}

\vfill
\newpage
\prob Compute the flux of the vector field
	\[\vec{F}(x,y,z) = (x^2+y^2) \vec{i} + (y^2 + z^2) \vec{j} + (x^2+z^2) \vec{k}
	\]
through the surface $S$ of the cube below (oriented outward).

\begin{flushright}
\includegraphics[scale=1.5]{pic1.pdf}
\end{flushright}

\vfill
\prob Compute the flux of $\vec{F}(x,y,z) = \ang{xy,yz,xz}$ through the sphere $S$ of radius $1$ centered at the origin (oriented outward).
\vfill

\newpage
\prob Let $S$ be an outward-oriented closed cylinder of height $2$ with its base a disk of radius $2$ in the $xy$-plane centered at the origin. Compute the flux of $\vec{F}(x,y,z) = \ang{0,y,0}$ through $S$ in two different ways.
\end{document}